\DocumentMetadata{
  pdfversion=2.0,
  pdfstandard=ua-2,
  lang=en-US,
  tagging=on,
  tagging-setup={math/setup=mathml-SE,tabsorder=structure}
}
\documentclass[11pt,letterpaper]{article}
\usepackage[margin=0.68in,includefoot]{geometry}
\usepackage{newcomputermodern}
\usepackage{amsmath,amscd,mathtools}
\usepackage{tikz,adjustbox}
\providecommand{\square}{\mdlgwhtsquare}
\usepackage[hidelinks]{hyperref}
\hypersetup{
  pdftitle={Algebra First-Year Exam, May 2022, Part 2},
  pdfauthor={University of Florida Department of Mathematics},
  pdfsubject={Graduate mathematics examination},
  pdfdisplaydoctitle=true,
  bookmarksopen=true
}
\setcounter{secnumdepth}{0}
\setlength{\parindent}{0pt}
\setlength{\parskip}{0.28em}
\setlength{\emergencystretch}{3em}
\setlength{\leftmargini}{2.15em}
\setlength{\leftmarginii}{2.65em}
\setlength{\leftmarginiii}{3em}
\renewcommand{\labelenumi}{\textbf{\theenumi.}}
\pagestyle{empty}
\raggedbottom
\allowdisplaybreaks
\newenvironment{examproblems}[1]{%
  \begin{enumerate}%
  \setcounter{enumi}{#1}%
  \setlength{\topsep}{0.35em}%
  \setlength{\partopsep}{0pt}%
  \setlength{\itemsep}{0.48em}%
  \setlength{\parsep}{0.18em}%
}{\end{enumerate}}
\newenvironment{examsubparts}{%
  \begin{enumerate}%
  \setlength{\topsep}{0.18em}%
  \setlength{\partopsep}{0pt}%
  \setlength{\itemsep}{0.16em}%
  \setlength{\parsep}{0.1em}%
}{\end{enumerate}}
\newenvironment{examinstructions}{%
  \par\begingroup\itshape
}{%
  \par\endgroup\vspace{0.18em}
}
\makeatletter
\renewcommand\section{\@startsection{section}{1}{0pt}%
  {-1.25ex plus -.25ex minus -.1ex}%
  {0.55ex plus .1ex}%
  {\normalfont\large\bfseries}}
\renewcommand{\@maketitle}{%
  \newpage\null\vskip -1em%
  {\footnotesize\bfseries
    \href{https://www.ufl.edu/}{University of Florida}, \href{https://math.ufl.edu/}{Department of Mathematics}\par}%
  \vskip -0.5em%
  \tagstructbegin{tag=Title}%
    {\LARGE\bfseries\href{https://gma.math.ufl.edu/exams/algebra-first-year/}{Algebra First-Year Exam}, May 2022, Part 2\par}%
  \tagstructend%
  \vskip 0.25em%
  \tagmcbegin{artifact}%
    \hrule height 1pt%
  \tagmcend%
  \vskip 1em%
}
\makeatother
\title{Algebra First-Year Exam, May 2022, Part 2}
\author{}
\date{}
\begin{document}
\maketitle
\thispagestyle{empty}
\begin{examinstructions}
Answer 4 questions. Write your solutions clearly in complete English sentences. You may quote results (within reason) as long as you state them clearly.
\end{examinstructions}
\begin{examproblems}{0}
\item
Prove that the following polynomials are irreducible in \(\mathbb{Q}[x]\), stating clearly any theorems you wish to apply.
\begin{examsubparts}
\item[\textnormal{(a)}]
\(x^4-75\).
\item[\textnormal{(b)}]
\(x^3-5x^2+x+5\).
\item[\textnormal{(c)}]
\(x^4-14x^3+2x^2+21x-7\).
\end{examsubparts}
\item
Let~\(R\) be the quotient of the ring \(\mathbb{Z}[i]\) of Gaussian integers by the principal ideal \((3+3i)\). Show that~\(R\) is finite and determine the number of its elements.
\item
Let~\(R\) be a commutative ring with 1 and let~\(I\), \(J\) and~\(P\) be ideals in~\(R\) such that \(I\cap J\subseteq P\).
\begin{examsubparts}
\item[\textnormal{(a)}]
Prove that if~\(P\) is a prime ideal then either \(I\subseteq P\) or \(J\subseteq P\).
\item[\textnormal{(b)}]
Give an example which shows that if~\(P\) is not prime then the conclusion above may not hold.
\end{examsubparts}
\item
Let~\(R\) be a ring with 1 and~\(M\) a (unital) left~\(R\)-module.
\begin{examsubparts}
\item[\textnormal{(a)}]
State what it means for~\(M\) to satisfy the Ascending Chain Condition (i.e. for~\(M\) to be Noetherian).
\item[\textnormal{(b)}]
Prove that if~\(M\) satisfies the ACC then every submodule of~\(M\) is finitely generated.
\item[\textnormal{(c)}]
Give an example of a finitely generated module that does not satisfy the ACC.
\end{examsubparts}
\item
Determine the conjugacy classes of elements of order 8 in \(\mathrm{GL}(4,\mathbb{Q}(\sqrt{2}))\). (Hint: Factorize \(x^4+1\) in \(\mathbb{Q}(\sqrt{2})[x]\).)
\item
Let~\(E\) be an extension field of the field~\(F\).
\begin{examsubparts}
\item[\textnormal{(a)}]
What does it mean for~\(E\) to be an algebraic extension of~\(F\)?
\item[\textnormal{(b)}]
Show that if~\(\alpha\) and~\(\beta\) are elements of~\(E\) that are both roots of the same irreducible polynomial in \(F[x]\), then the subfields \(F(\alpha)\) and \(F(\beta)\) of~\(E\) are isomorphic.
\item[\textnormal{(c)}]
Give an example where the subfields in (b) are not equal.
\end{examsubparts}
\end{examproblems}
\end{document}
